\PassOptionsToPackage{usenames,dvipsnames}{xcolor}
\documentclass[a4paper,12pt]{letter}
\usepackage[utf8]{inputenc}
\usepackage[T1]{fontenc}
\usepackage[english,slovene]{babel}
\usepackage[tikz]{bclogo}
\usepackage{qrcode,marvosym}
\usepackage{pgf,tikz,pgfplots}
\pgfplotsset{compat=1.14}
\usepackage{mathrsfs}
\usetikzlibrary{arrows,angles,quotes}
\usepackage{graphicx}
\newcommand{\ds}{\displaystyle}
\usepackage{fancybox}
\usepackage{amsfonts}
\usepackage{amsmath}
\usepackage{wallpaper}
\usepackage{hyperref}
\usepackage{array}
\newcolumntype{M}[1]{>{\centering\arraybackslash}m{#1}}
\newcolumntype{N}{@{}m{0pt}@{}}
\usepackage{fancyhdr,enumerate}
\newcommand{\ora}{\overrightarrow}
\renewcommand{\headrulewidth}{3pt}
\renewcommand{\headrule}{\hbox to\headwidth{%
  \color{gray}\leaders\hrule height \headrulewidth\hfill}}
\renewcommand{\footrulewidth}{2pt}
\renewcommand{\footrule}{\hbox to\headwidth{%
  \color{brown}\leaders\hrule height \footrulewidth\hfill}}
\usepackage{gensymb}
\usepackage{amssymb}
  \textwidth 20 true cm
  \oddsidemargin -2 true cm
  \evensidemargin -2 true cm
  \topmargin -3.5 true cm
  \textheight 27.5 true cm
  \linespread{1.2}
\renewcommand{\headrulewidth}{0.3pt}
\renewcommand{\footrulewidth}{1.9pt}
\definecolor{qqwuqq}{rgb}{1,0,0}
\definecolor{qqqqff}{rgb}{0,0,1}
\definecolor{qqffqq}{rgb}{0,1,0}
\definecolor{ffqqqq}{rgb}{1,0,0}
\definecolor{zzuxux}{rgb}{0,0,0}
\usepackage{mdframed}
\usepackage{multicol}
\definecolor{bgblue}{RGB}{0,0,253}
\usepackage{eurosym}
%%%%%%%%%%%%%%%%%%%%%%%%%%%%%%%%%%%%%%%%%%%%%%%%%%%%%%%%
\newcounter{tocke}
\setcounter{tocke}{0}
\usepackage{tikz-3dplot}
%%%%%%%%%%%%%%%%%%%%%%%%%%%%%%%%%%%%%%%%%%%%%%%%%%%%%%%%%%
\mdfdefinestyle{theoremstyle}{%
linecolor=brown!40,linewidth=1pt,%
frametitlerule=true,%
frametitlebackgroundcolor=brown!50,
innertopmargin=\topskip,
}
\mdfdefinestyle{theoremstyle1}{%
linecolor=brown,middlelinewidth=1pt,%
frametitlerule=true,%
apptotikzsetting={\tikzset{mdfframetitlebackground/.append style={%
shade,left color=black!40, right color=gray!10},
mdfbackground/.append style={
top color=white!100!white,
bottom color=black!5
},
}},
frametitlerulecolor=gray,
frametitlerulewidth=2pt,
innertopmargin=\topskip,
innerbottommargin=\topskip,
}
\mdtheorem[style=theoremstyle1]{naloga}{Naloga}
\mdtheorem[style=theoremstyle]{resitev}{Analiza Naloge}
\mdtheorem[style=theoremstyle]{kriterij}{\v{S}tevilo dose\v{z}enih to\v{c}k na testu}
%%%%%%%%%%%%%%%%%%%%%%%%%%%%%%%%%%%%%%%%%%%%%%%%%%%%%%%%%%
\usepackage[table]{xcolor}
\usepackage[printwatermark]{xwatermark}
\newwatermark[allpages,color=brown!2,angle=45,scale=5,
xpos=-5,ypos=0]{matej.info}
%%%%%%%%%%%%%%%%%%%%%%%%%%%%%%%%%%%%%%%%%%%%%%%%%%%%%%%%%%

\begin{document}
\pagestyle{fancy}
\renewcommand\headrulewidth{0pt}
\lhead{}\chead{}\rhead{}
\rfoot{\vspace*{-2\baselineskip}}
\renewcommand{\vec}{\overrightarrow}
%%%%%%%%%%%%%%%%%%%% podatki testa -- tu spreminjaj
\def\programletnik{G-3}
\def\naslov{Preverjanje 2.4 - }
\def\sklop{Kotne funkcije}
%%%%%%%%%%%%%%%%%%%% tocke po nalogah
\def\tockeA{1+1+1+1+1+1}
\def\tockeB{2+3}
\def\tockeC{2+3}
\def\tockeD{3+3}
\def\tockeE{3}
\def\tockeF{3+3}
%%%%%%%%%%%%%%%%%%%%
\begin{minipage}{20cm}
\begin{bclogo}[couleur=brown!30,
  arrondi=0,ombre=true, couleurOmbre=gray!50,
logo =\bcplume,
  noborder=true]
{\textrm{{\Large{\naslov \sklop} \Huge{}
\large{} \hfill $\programletnik$}}}
\end{bclogo}
\end{minipage}
\begin{minipage}{20cm}
\begin{bclogo}[couleur=brown!30,
  arrondi=0,ombre=true, couleurOmbre=gray!50,
logo =\bcplume,
  noborder=true]
{\textrm{{\raisebox{-0.3cm}{\large{Ime in Priimek: \rule{10cm}{0.1mm}} \Huge{}
\large{} \hfill $ $}}}}
\end{bclogo}
\end{minipage}

%%%%%%%%%%%%%%%%%%%%%%%%%%%%%%%%%%%%%%%%%%%%%%%%% (po zgledu 2N.2 -- 1i)
\begin{naloga}[\hfill $\tockeA$] \addtocounter{tocke}{\numexpr\tockeA\relax}
Dano je $\sin x=\frac45$, $x\in\left(\frac{\pi}{2},\pi\right)$, in $\cos y=-\frac{5}{13}$, $y\in\left(\pi,\frac{3\pi}{2}\right)$. Brez kalkulatorja izra\v{c}unaj (vrednosti zapi\v{s}i natan\v{c}no):
\begin{multicols}{3}
\begin{enumerate}[a)]
\item $\ds\sin(x+y)$
\item $\ds\cos(x-y)$
\item $\ds\tan\left(x+\frac{\pi}{4}\right)$
\item $\ds\cos2y$
\item $\ds\sin\frac{x}{2}$
\item $\ds\tan\frac{y}{2}$
\end{enumerate}
\end{multicols}
\end{naloga}\newpage

%%%%%%%%%%%%%%%%%%%%%%%%%%%%%%%%%%%%%%%%%%%%%%%%% (po zgledu 2N.2 -- 4a, 4g)
\begin{naloga}[\hfill $\tockeB$] \addtocounter{tocke}{\numexpr\tockeB\relax}
Doka\v{z}i trigonometri\v{c}ni enakosti.
\begin{enumerate}[a)]
\item $\ds(\cos\alpha-\sin\alpha)^2=1-\sin2\alpha$
\item $\ds\frac{\cos\alpha+\sin\alpha}{\cos\alpha-\sin\alpha}=\frac{1+\sin2\alpha}{\cos2\alpha}$
\end{enumerate}
\end{naloga}\newpage

%%%%%%%%%%%%%%%%%%%%%%%%%%%%%%%%%%%%%%%%%%%%%%%%% (po zgledu 2N.2 -- 6, 2N.3 -- 1e)
\begin{naloga}[\hfill $\tockeC$] \addtocounter{tocke}{\numexpr\tockeC\relax}
Poenostavi izraza.
\begin{enumerate}[a)]
\item $\ds\sin\left(x+\frac{\pi}{4}\right)+\cos\left(x+\frac{\pi}{4}\right)$
\item $\ds\frac{\sin2x+\sin4x+\sin6x}{\cos2x+\cos4x+\cos6x}$
\end{enumerate}
\end{naloga}\newpage

%%%%%%%%%%%%%%%%%%%%%%%%%%%%%%%%%%%%%%%%%%%%%%%%% (po zgledu 2N.3 -- 2c)
\begin{naloga}[\hfill $\tockeD$] \addtocounter{tocke}{\numexpr\tockeD\relax}
Nari\v{s}i graf funkcije $\ds f(x)=2\sin\Big(x-\frac{\pi}{3}\Big)+1$ s pomo\v{c}jo transformacij osnovnega grafa $y=\sin x$ (premiki, raztegi). Vsak vmesni korak nari\v{s}i v isti koordinatni sistem. Na intervalu $[-\pi,2\pi]$ dolo\v{c}i zalogo vrednosti, ni\v{c}le in ekstreme ter izra\v{c}unaj za\v{c}etno vrednost $f(0)$.
\begin{center}
\begin{tikzpicture}[xscale=1.6,yscale=1.1,>=stealth]
\draw[very thin,gray!30,xstep=0.785398,ystep=0.5] (-3.1416,-1.5) grid (6.2832,3.5);
\draw[->,thick] (-3.4416,0) -- (6.7832,0) node[below left] {$x$};
\draw[->,thick] (0,-1.8) -- (0,4) node[below left] {$y$};
\draw (-3.1416,0.1143) -- (-3.1416,-0.1143);
\node[below,inner sep=2pt] at (-3.1416,-0.1429) {\scriptsize $-\pi$};
\draw (-1.5708,0.1143) -- (-1.5708,-0.1143);
\node[below,inner sep=2pt] at (-1.5708,-0.1429) {\scriptsize $-\tfrac{\pi}{2}$};
\draw (1.5708,0.1143) -- (1.5708,-0.1143);
\node[below,inner sep=2pt] at (1.5708,-0.1429) {\scriptsize $\tfrac{\pi}{2}$};
\draw (3.1416,0.1143) -- (3.1416,-0.1143);
\node[below,inner sep=2pt] at (3.1416,-0.1429) {\scriptsize $\pi$};
\draw (4.7124,0.1143) -- (4.7124,-0.1143);
\node[below,inner sep=2pt] at (4.7124,-0.1429) {\scriptsize $\tfrac{3\pi}{2}$};
\draw (6.2832,0.1143) -- (6.2832,-0.1143);
\node[below,inner sep=2pt] at (6.2832,-0.1429) {\scriptsize $2\pi$};
\draw (0.0800,-1) -- (-0.0800,-1) node[left,inner sep=2pt] {\scriptsize $-1$};
\draw (0.0800,1) -- (-0.0800,1) node[left,inner sep=2pt] {\scriptsize $1$};
\draw (0.0800,2) -- (-0.0800,2) node[left,inner sep=2pt] {\scriptsize $2$};
\draw (0.0800,3) -- (-0.0800,3) node[left,inner sep=2pt] {\scriptsize $3$};
\node[below left,inner sep=2pt] at (0,0) {\scriptsize $0$};
\end{tikzpicture}
\end{center}
\begin{center}$Z_f=$\,\rule{4cm}{0.1mm}\qquad $f(0)=$\,\rule{4cm}{0.1mm}\end{center}
\vspace{0.3cm}
Ni\v{c}le: \rule{12cm}{0.1mm}\\[0.5cm]
Ekstremi: \rule{11.6cm}{0.1mm}
\end{naloga}\newpage

%%%%%%%%%%%%%%%%%%%%%%%%%%%%%%%%%%%%%%%%%%%%%%%%% (po zgledu 2N.3 -- 3)
\begin{naloga}[\hfill $\tockeE$] \addtocounter{tocke}{\numexpr\tockeE\relax}
Iz grafa razberi funkcijski predpis funkcije $f$ oblike $f(x)=A\sin(bx+c)$.
\begin{center}
\begin{tikzpicture}[xscale=1.4,yscale=0.8,>=stealth]
\draw[very thin,gray!30,xstep=0.523599,ystep=0.5] (-3.1416,-2.5) grid (6.2832,2.5);
\draw[->,thick] (-3.4416,0) -- (6.7832,0) node[below left] {$x$};
\draw[->,thick] (0,-2.8) -- (0,3) node[below left] {$y$};
\draw (-3.1416,0.1143) -- (-3.1416,-0.1143);
\node[below,inner sep=2pt] at (-3.1416,-0.1429) {\scriptsize $-\pi$};
\draw (-1.5708,0.1143) -- (-1.5708,-0.1143);
\node[below,inner sep=2pt] at (-1.5708,-0.1429) {\scriptsize $-\tfrac{\pi}{2}$};
\draw (1.5708,0.1143) -- (1.5708,-0.1143);
\node[below,inner sep=2pt] at (1.5708,-0.1429) {\scriptsize $\tfrac{\pi}{2}$};
\draw (3.1416,0.1143) -- (3.1416,-0.1143);
\node[below,inner sep=2pt] at (3.1416,-0.1429) {\scriptsize $\pi$};
\draw (4.7124,0.1143) -- (4.7124,-0.1143);
\node[below,inner sep=2pt] at (4.7124,-0.1429) {\scriptsize $\tfrac{3\pi}{2}$};
\draw (6.2832,0.1143) -- (6.2832,-0.1143);
\node[below,inner sep=2pt] at (6.2832,-0.1429) {\scriptsize $2\pi$};
\draw (0.0800,-2) -- (-0.0800,-2) node[left,inner sep=2pt] {\scriptsize $-2$};
\draw (0.0800,-1) -- (-0.0800,-1) node[left,inner sep=2pt] {\scriptsize $-1$};
\draw (0.0800,1) -- (-0.0800,1) node[left,inner sep=2pt] {\scriptsize $1$};
\draw (0.0800,2) -- (-0.0800,2) node[left,inner sep=2pt] {\scriptsize $2$};
\node[below left,inner sep=2pt] at (0,0) {\scriptsize $0$};
\begin{scope}
\clip (-3.1416,-2.5) rectangle (6.2832,2.5);
\draw[red,line width=1.6pt,domain=-3.1416:6.2832,samples=90,smooth] plot (\x,{2*sin(deg(2*\x-pi/3))});
\end{scope}
\end{tikzpicture}
\end{center}
\begin{center}$f(x)=$\,\rule{6cm}{0.1mm}\end{center}
\end{naloga}\newpage

%%%%%%%%%%%%%%%%%%%%%%%%%%%%%%%%%%%%%%%%%%%%%%%%%
\begin{naloga}[\hfill $\tockeF$] \addtocounter{tocke}{\numexpr\tockeF\relax}
Brez kalkulatorja izra\v{c}unaj natan\v{c}no vrednost izraza (upo\v{s}tevaj periodi\v{c}nost, sodost oziroma lihost in prehod na ostri kot).
\begin{enumerate}[a)]
\item $\ds\frac{\sin750\degree\cdot\cos(-300\degree)+\tan1125\degree}{\cos840\degree-\sin(-210\degree)}$
\item $\ds\frac{\sin\dfrac{17\pi}{6}+\cos\left(-\dfrac{11\pi}{3}\right)}{\tan\dfrac{13\pi}{4}\cdot\cos\dfrac{7\pi}{6}}$
\end{enumerate}
\end{naloga}

\end{document}

\vfill
\begin{kriterij*} \begin{center} \renewcommand{\arraystretch}{2.5}
\begin{tabular}{|>{\columncolor{gray!20}}c||c|c|c|c|c||c|c|c|} \hline Odstotki &$[0,45)$ & $[45,60)$ & $[60,75)$ & $[75,90)$ & $[90,100]$& \v{S}t. mo\v{z}nih to\v{c}k& Dose\v{z}ene to\v{c}ke & OCENA \\ \hline Ocena & 1 & 2 & 3 & 4 & 5 &\cellcolor{gray!20} $\thetocke$&\cellcolor{gray!20} \phantom{ocena} & \cellcolor{gray!20} \phantom{a} \\ \hline \end{tabular} \end{center} \vspace{0.3cm}  \end{kriterij*}

\newpage
\textsc{\Large Re\v{s}itve s postopki}
\vspace{0.4cm}

%%%%%%%%%%%%%%%%%%%% pregled rezultatov
\begin{center}
\small
\renewcommand{\arraystretch}{2.1}
\begin{tabular}{|>{\columncolor{gray!20}}c|l|}
\hline
\textbf{Naloga} & \textbf{Rezultati} \\ \hline
\textbf{1} & a) $\frac{16}{65}$\quad b) $-\frac{33}{65}$\quad c) $-\frac17$\quad d) $-\frac{119}{169}$\quad e) $\frac{2\sqrt5}{5}$\quad f) $-\frac32$ \\ \hline
\textbf{2} & dokaz enakosti \\ \hline
\textbf{3} & a) $\sqrt2\cos x$\quad b) $\tan4x$ \\ \hline
\textbf{4} & $Z_f=[-1,3]$; ni\v{c}le $-\frac{\pi}{2},\frac{\pi}{6},\frac{3\pi}{2}$; $\max\,(\frac{5\pi}{6},3)$, $\min\,(-\frac{\pi}{6},-1),(\frac{11\pi}{6},-1)$; $f(0)=1-\sqrt3$ \\ \hline
\textbf{5} & $f(x)=2\sin\left(2x-\frac{\pi}{3}\right)$ \\ \hline
\textbf{6} & a) $-\frac54$\quad b) $-\frac{2\sqrt3}{3}$ \\ \hline
\end{tabular}
\end{center}
\vspace{0.2cm}

%%%%%%%%%%%%%%%%%%%% naloga 1
\begin{resitev}[-- adicijski izreki, dvojni in polovi\v{c}ni koti]
Manjkajo\v{c}e vrednosti: $x$ le\v{z}i v II. kvadrantu, $y$ v III. kvadrantu:
\[
\cos x=-\sqrt{1-\tfrac{16}{25}}=-\frac35,\quad \tan x=-\frac43,\qquad \sin y=-\sqrt{1-\tfrac{25}{169}}=-\frac{12}{13},\quad \tan y=\frac{12}{5}.
\]
\textbf{a)} $\ds\sin(x+y)=\sin x\cos y+\cos x\sin y=\frac45\cdot\Big(-\frac{5}{13}\Big)+\Big(-\frac35\Big)\Big(-\frac{12}{13}\Big)=\frac{-20+36}{65}=\boxed{\frac{16}{65}}$

\textbf{b)} $\ds\cos(x-y)=\cos x\cos y+\sin x\sin y=\Big(-\frac35\Big)\Big(-\frac{5}{13}\Big)+\frac45\cdot\Big(-\frac{12}{13}\Big)=\frac{15-48}{65}=\boxed{-\frac{33}{65}}$

\textbf{c)} $\ds\tan\Big(x+\frac{\pi}{4}\Big)=\frac{\tan x+1}{1-\tan x}=\frac{-\frac43+1}{1+\frac43}=\frac{-\frac13}{\frac73}=\boxed{-\frac17}$

\textbf{d)} $\ds\cos2y=2\cos^2y-1=2\cdot\frac{25}{169}-1=\boxed{-\frac{119}{169}}$

\textbf{e)} $\frac{x}{2}\in\left(\frac{\pi}{4},\frac{\pi}{2}\right)$, zato je $\sin\frac x2>0$:
$\ds\sin\frac{x}{2}=\sqrt{\frac{1-\cos x}{2}}=\sqrt{\frac{1+\frac35}{2}}=\sqrt{\frac45}=\boxed{\frac{2\sqrt5}{5}}$

\textbf{f)} $\frac{y}{2}\in\left(\frac{\pi}{2},\frac{3\pi}{4}\right)$, zato je $\tan\frac y2<0$:
$\ds\tan\frac{y}{2}=\frac{\sin y}{1+\cos y}=\frac{-\frac{12}{13}}{1-\frac{5}{13}}=\frac{-\frac{12}{13}}{\frac{8}{13}}=\boxed{-\frac32}$
\end{resitev}

%%%%%%%%%%%%%%%%%%%% naloga 2
\begin{resitev}[-- dokazi enakosti]
\textbf{a)}
\[
(\cos\alpha-\sin\alpha)^2=\cos^2\alpha-2\sin\alpha\cos\alpha+\sin^2\alpha=1-\sin2\alpha.
\]
\textbf{b)} Ulomek raz\v{s}irimo s $\cos\alpha+\sin\alpha$:
\[
\frac{\cos\alpha+\sin\alpha}{\cos\alpha-\sin\alpha}=\frac{(\cos\alpha+\sin\alpha)^2}{\cos^2\alpha-\sin^2\alpha}=\frac{\cos^2\alpha+2\sin\alpha\cos\alpha+\sin^2\alpha}{\cos2\alpha}=\frac{1+\sin2\alpha}{\cos2\alpha}.
\]
\end{resitev}

%%%%%%%%%%%%%%%%%%%% naloga 3
\begin{resitev}[-- poenostavljanje izrazov]
\textbf{a)}
\[
\sin x\cdot\frac{\sqrt2}{2}+\cos x\cdot\frac{\sqrt2}{2}+\cos x\cdot\frac{\sqrt2}{2}-\sin x\cdot\frac{\sqrt2}{2}=\boxed{\sqrt2\cos x}
\]
\textbf{b)} Zdru\v{z}imo prvi in tretji \v{c}len (faktorizacija):
\[
\sin2x+\sin6x=2\sin4x\cos2x,\qquad \cos2x+\cos6x=2\cos4x\cos2x.
\]
\[
\frac{2\sin4x\cos2x+\sin4x}{2\cos4x\cos2x+\cos4x}=\frac{\sin4x\,(2\cos2x+1)}{\cos4x\,(2\cos2x+1)}=\boxed{\tan4x}
\]
\end{resitev}

%%%%%%%%%%%%%%%%%%%% naloga 4
\begin{resitev}[-- graf, zaloga vrednosti, ni\v{c}le, ekstremi]
Koraki: $y=\sin x\ \xrightarrow{\text{premik desno za }\frac{\pi}{3}}\ y=\sin\big(x-\frac{\pi}{3}\big)\ \xrightarrow{\text{raztezek }\times2}\ y=2\sin\big(x-\frac{\pi}{3}\big)\ \xrightarrow{\text{premik gor za }1}\ f(x)$.
\begin{center}
\begin{tikzpicture}[xscale=1.5,yscale=1,>=stealth]
\draw[very thin,gray!30,xstep=0.785398,ystep=0.5] (-3.1416,-1.5) grid (6.2832,3.5);
\draw[->,thick] (-3.4416,0) -- (6.7832,0) node[below left] {$x$};
\draw[->,thick] (0,-1.8) -- (0,4) node[below left] {$y$};
\draw (-3.1416,0.1143) -- (-3.1416,-0.1143);
\node[below,inner sep=2pt] at (-3.1416,-0.1429) {\scriptsize $-\pi$};
\draw (-1.5708,0.1143) -- (-1.5708,-0.1143);
\node[below,inner sep=2pt] at (-1.5708,-0.1429) {\scriptsize $-\tfrac{\pi}{2}$};
\draw (1.5708,0.1143) -- (1.5708,-0.1143);
\node[below,inner sep=2pt] at (1.5708,-0.1429) {\scriptsize $\tfrac{\pi}{2}$};
\draw (3.1416,0.1143) -- (3.1416,-0.1143);
\node[below,inner sep=2pt] at (3.1416,-0.1429) {\scriptsize $\pi$};
\draw (4.7124,0.1143) -- (4.7124,-0.1143);
\node[below,inner sep=2pt] at (4.7124,-0.1429) {\scriptsize $\tfrac{3\pi}{2}$};
\draw (6.2832,0.1143) -- (6.2832,-0.1143);
\node[below,inner sep=2pt] at (6.2832,-0.1429) {\scriptsize $2\pi$};
\draw (0.0800,-1) -- (-0.0800,-1) node[left,inner sep=2pt] {\scriptsize $-1$};
\draw (0.0800,1) -- (-0.0800,1) node[left,inner sep=2pt] {\scriptsize $1$};
\draw (0.0800,2) -- (-0.0800,2) node[left,inner sep=2pt] {\scriptsize $2$};
\draw (0.0800,3) -- (-0.0800,3) node[left,inner sep=2pt] {\scriptsize $3$};
\node[below left,inner sep=2pt] at (0,0) {\scriptsize $0$};
\begin{scope}
\clip (-3.1416,-1.5) rectangle (6.2832,3.5);
\draw[gray,dashed,thick,domain=-3.1416:6.2832,samples=60,smooth] plot (\x,{sin(deg(\x))});
\end{scope}
\begin{scope}
\clip (-3.1416,-1.5) rectangle (6.2832,3.5);
\draw[blue,dashed,thick,domain=-3.1416:6.2832,samples=60,smooth] plot (\x,{sin(deg(\x-pi/3))});
\end{scope}
\begin{scope}
\clip (-3.1416,-1.5) rectangle (6.2832,3.5);
\draw[green!60!black,dashed,thick,domain=-3.1416:6.2832,samples=60,smooth] plot (\x,{2*sin(deg(\x-pi/3))});
\end{scope}
\begin{scope}
\clip (-3.1416,-1.5) rectangle (6.2832,3.5);
\draw[red,line width=1.6pt,domain=-3.1416:6.2832,samples=90,smooth] plot (\x,{2*sin(deg(\x-pi/3))+1});
\end{scope}
\foreach \p in {(-1.5708,0),(0.5236,0),(4.7124,0),(2.6180,3),(-0.5236,-1),(5.7596,-1),(0,-0.7321)} \fill[red] \p circle (2.2pt);
\node[red,right] at (5.2,3.2) {\scriptsize $f(x)$};
\end{tikzpicture}
\end{center}
\begin{itemize}
\item Zaloga vrednosti: $\sin\big(x-\frac{\pi}{3}\big)\in[-1,1]\Rightarrow f(x)\in[-1,3]$, \quad $\boxed{Z_f=[-1,3]}$
\item Ni\v{c}le: $2\sin\big(x-\frac{\pi}{3}\big)+1=0\Rightarrow\sin\big(x-\frac{\pi}{3}\big)=-\frac12\Rightarrow x-\frac{\pi}{3}=-\frac{\pi}{6}+2k\pi$ ali $x-\frac{\pi}{3}=\frac{7\pi}{6}+2k\pi$, torej $x=\frac{\pi}{6}+2k\pi$ ali $x=\frac{3\pi}{2}+2k\pi$. Na $[-\pi,2\pi]$: $\boxed{x_1=-\frac{\pi}{2},\ x_2=\frac{\pi}{6},\ x_3=\frac{3\pi}{2}}$
\item Maksimum: $x-\frac{\pi}{3}=\frac{\pi}{2}\Rightarrow x=\frac{5\pi}{6}$, \quad $\boxed{T_{\max}\big(\frac{5\pi}{6},3\big)}$
\item Minimum: $x-\frac{\pi}{3}=-\frac{\pi}{2}+2k\pi\Rightarrow x=-\frac{\pi}{6}+2k\pi$, \quad $\boxed{T_{\min}\big(-\frac{\pi}{6},-1\big),\ T_{\min}\big(\frac{11\pi}{6},-1\big)}$
\item Za\v{c}etna vrednost: $f(0)=2\sin\big(-\frac{\pi}{3}\big)+1=-\sqrt3+1$, \quad $\boxed{f(0)=1-\sqrt3\approx-0{,}73}$
\end{itemize}
\end{resitev}

\newpage
%%%%%%%%%%%%%%%%%%%% naloga 5
\begin{resitev}[-- predpis iz grafa]
\begin{itemize}
\item Amplituda: najve\v{c}ja vrednost $2$, najmanj\v{s}a $-2$, zato $A=2$.
\item Perioda: razdalja med sosednjima ni\v{c}lama je $\frac{\pi}{2}$, zato je osnovna perioda $\omega=\pi$ in $b=\frac{2\pi}{\omega}=2$.
\item Premik: graf nara\v{s}\v{c}ajo\v{c}e seka os $x$ pri $x=\frac{\pi}{6}$, zato $2\cdot\frac{\pi}{6}+c=0\Rightarrow c=-\frac{\pi}{3}$.
\end{itemize}
\[
\boxed{f(x)=2\sin\Big(2x-\frac{\pi}{3}\Big)=2\sin2\Big(x-\frac{\pi}{6}\Big)}
\]
\end{resitev}

%%%%%%%%%%%%%%%%%%%% naloga 6
\begin{resitev}[-- periodi\v{c}nost, lihost, prehod na ostri kot]
\textbf{a)}
\[
\sin750\degree=\sin(30\degree+2\cdot360\degree)=\sin30\degree=\tfrac12,\qquad \cos(-300\degree)=\cos300\degree=\cos60\degree=\tfrac12,
\]
\[
\tan1125\degree=\tan(45\degree+6\cdot180\degree)=1,\quad \cos840\degree=\cos120\degree=-\cos60\degree=-\tfrac12,\quad \sin(-210\degree)=-\sin210\degree=\sin30\degree=\tfrac12.
\]
\[
\frac{\frac12\cdot\frac12+1}{-\frac12-\frac12}=\frac{\frac54}{-1}=\boxed{-\frac54}
\]
\textbf{b)}
\[
\sin\frac{17\pi}{6}=\sin\frac{5\pi}{6}=\sin\frac{\pi}{6}=\tfrac12,\qquad \cos\Big(-\frac{11\pi}{3}\Big)=\cos\frac{11\pi}{3}=\cos\Big(4\pi-\frac{\pi}{3}\Big)=\cos\frac{\pi}{3}=\tfrac12,
\]
\[
\tan\frac{13\pi}{4}=\tan\Big(3\pi+\frac{\pi}{4}\Big)=1,\qquad \cos\frac{7\pi}{6}=\cos\Big(\pi+\frac{\pi}{6}\Big)=-\cos\frac{\pi}{6}=-\frac{\sqrt3}{2}.
\]
\[
\frac{\frac12+\frac12}{1\cdot\left(-\frac{\sqrt3}{2}\right)}=-\frac{2}{\sqrt3}=\boxed{-\frac{2\sqrt3}{3}}
\]
\end{resitev}

\end{document}
